% Theory section. All numbers are \input from generated/ (scripts/make_tables.py).
\section{Theory: when is hindsight planning better, and by how much?}
\label{sec:theory}

\paragraph{Setting.} At step $t$ the agent receives a commanded displacement $\delta_t$ (i.i.d., independent of
the past) and chooses a bounded command $|a_t|\le A$. The world responds
\begin{equation}
\Delta s_t=\theta_t a_t+w_t+\varepsilon_t,\qquad \varepsilon_t\sim\mathcal N(0,\sigma^2)\ \text{i.i.d.},
\end{equation}
with cost $c_t=(\Delta s_t-\delta_t)^2$. The hidden disturbance $\phi_t=(\theta_t,w_t)$ has a \emph{gain} part
$\theta_t$ (it scales the command) and an \emph{additive} part $w_t$; its prior is $\mathcal N(m_0,S_0)$ with
$m_0=(\bar\theta,0)$. Let $\mathcal F_t$ be the history $(\delta_{0:t},a_{0:t-1},\Delta s_{0:t-1})$.
A \textbf{forward} planner chooses $a_t$ as a function of $\delta_t$ only (it never conditions on outcomes);
a \textbf{hindsight} planner chooses $a_t$ as a function of $\mathcal F_t$. Regret is measured against the
oracle that knows $\phi_t$ and obeys the same bound; $\varepsilon_t$ is common to all planners and cancels.

\begin{theorem}[Dominance]\label{thm:dominance}
For any prior and any disturbance process, the myopic Bayes hindsight planner
$a^H_t=\arg\min_{|a|\le A}\mathbb E[c_t\mid\mathcal F_t]$, computed under the correct disturbance model,
satisfies $\mathbb E c^H_t\le\mathbb E c^F_t$ for every $t$ and every forward planner $F$.
\end{theorem}
\begin{proof}
$a^F_t$ is a function of $\delta_t$, hence $\mathcal F_t$-measurable, so it is feasible for the minimisation defining
$a^H_t$: $\mathbb E[c^H_t\mid\mathcal F_t]\le\mathbb E[c(a^F_t)\mid\mathcal F_t]$. Because $\delta_t$ is
independent of the past and actions do not move the state that later costs depend on, the law of $c(a^F_t)$
under the hindsight planner's history equals that under $F$'s own history. Take expectations.
This is Blackwell's comparison of experiments~\citep{blackwell1953} specialised to one-step costs.
\end{proof}

\begin{theorem}[No free lunch under fresh disturbances]\label{thm:fresh}
If $\phi_t$ are i.i.d.\ and independent of the past, the posterior predictive of $\phi_t$ given $\mathcal F_t$ equals
the prior, so the correctly specified hindsight planner coincides with the Bayes forward planner and
$\mathbb E c^H_t=\mathbb E c^F_t$. A hindsight planner that \emph{assumes} persistence is misspecified and can be
strictly worse (Table~\ref{tab:theory}, row \emph{fresh}).
\end{theorem}

\begin{theorem}[Linear vs.\ logarithmic regret under persistent disturbances]\label{thm:separation}
Let $\phi_t\equiv\phi\sim\mathcal N(m_0,S_0)$ be persistent and assume targets are reachable
($|\delta_t-w|\le A|\theta|$). Then
\begin{enumerate}
\item the Bayes forward planner pays a constant excess every step,
  $c_F=\sigma_w^2+\sigma_\theta^2\,\mathbb E[\delta^2]/(\bar\theta^2+\sigma_\theta^2)$, so $R_F(T)=c_F\,T$;
\item the certainty-equivalent hindsight planner, and therefore the Bayes hindsight planner, satisfies
\begin{equation}
R_H(T)\le 2\sigma^2\max(1,\kappa)\,d\,\log\!\Big(\tfrac{\lambda_{\max}(V_0)+TL^2/d}{\lambda_{\min}(V_0)}\Big),\quad
V_0=\sigma^2S_0^{-1},\ L^2=A^2+1,\ \kappa=L^2/\lambda_{\min}(V_0),
\end{equation}
with $d=2$; i.e.\ $R_H(T)=O(\sigma^2 d\log T)$, and the per-step excess decays as $\sigma^2 d/t$.
\end{enumerate}
\end{theorem}
\begin{proof}[Proof sketch]
(1) Minimise $\mathbb E[(\theta a+w-\delta)^2]=(\bar\theta^2+\sigma_\theta^2)a^2-2a\bar\theta\delta+\delta^2+\sigma_w^2$
over $a$. (2) With $x_t=(a_t,1)$ and Gaussian conjugacy, the posterior is $\mathcal N(\mu_t,\Sigma_t)$ with
$\Sigma_t=\sigma^2V_t^{-1}$, $V_t=V_0+\sum_{s<t}x_sx_s^\top$. The CE action solves $x_t^\top\mu_t=\delta_t$, so its
error is $x_t^\top(\phi-\mu_t)$ and $\mathbb E[\,\cdot^2\mid\mathcal F_t]=x_t^\top\Sigma_tx_t=\sigma^2\|x_t\|^2_{V_t^{-1}}$.
Since $\|x_t\|^2_{V_t^{-1}}\le\kappa$, we have $u\le\max(1,\kappa)\min(1,u)$, and the elliptical-potential lemma
\citep[Lemma 11]{abbasi2011} gives $\sum_t\min(1,\|x_t\|^2_{V_{t}^{-1}})\le2\log(\det V_T/\det V_0)$; bound
$\det V_T$ by $(\lambda_{\max}(V_0)+TL^2/d)^d$. The Bayes planner is no worse than CE in posterior expectation at
every step because the CE action is feasible for its minimisation.
\end{proof}

\begin{theorem}[Residual reuse is inconsistent for sequential replanning]\label{thm:residual}
Apply the residual-reuse rule of \citet{cwm_2611_24778} at every step: predict
$\Delta s_t(a)=\bar\theta a+r_{t}$ with $r_t=\Delta s_{t-1}-\bar\theta a_{t-1}$, and act on it. Under persistent $\phi$
(unclipped), the error is $e_t=(\theta-\bar\theta)(a_t-a_{t-1})+\varepsilon_t-\varepsilon_{t-1}$. Hence
(i) with a pure additive fault ($\theta=\bar\theta$) it pays excess $\sigma^2$ \emph{every} step, $R(T)=\sigma^2T$,
versus $O(\log T)$ for hindsight; (ii) the commands follow an AR(1) recursion with coefficient
$-(\theta-\bar\theta)/\bar\theta$, which is unstable for $\theta>2\bar\theta$ or $\theta<0$;
(iii) with no disturbance at all it still pays $\sigma^2$ per step while the forward planner pays nothing.
\end{theorem}
\begin{proof}
Substitute $\bar\theta a_t=\delta_t-r_t$ into $e_t=\theta a_t+w+\varepsilon_t-\delta_t$.
\end{proof}
This does not contradict \citet{cwm_2611_24778}: their guarantee is for the counterfactual of \emph{the same event}
(same $\varepsilon$), which is exact under additivity. Item (iii) matches their own report that residual correction
lowers replanning success when there is no disturbance.

\paragraph{Numerical verification.} Table~\ref{tab:theory} checks each statement by Monte-Carlo
(\texttt{wmi/theory.py}, \texttt{scripts/run\_theory.py}).
\begin{table}[h]\centering\scriptsize
\IfFileExists{generated/theory_table.tex}{% generated by scripts/make_tables.py; seeds=[0, 1, 2]; commits=['6e7a3ea']; dirty=False
\begin{tabular}{lccccc}
\toprule
Regime & forward & residual & hindsight\_bayes & hindsight\_forget & $t\cdot$excess (hindsight) \\
\midrule
persistent & 1536.5 $\pm$ 15.2 & 13393.3 $\pm$ 2592.3 & 2.3 $\pm$ 0.0 & 5.3 $\pm$ 0.0 & 0.18 $\pm$ 0.00 \\
persistent\_additive & 1078.0 $\pm$ 11.2 & 360.1 $\pm$ 0.3 & 1.8 $\pm$ 0.0 & 4.8 $\pm$ 0.0 & 0.18 $\pm$ 0.00 \\
persistent\_gain & 543.9 $\pm$ 6.0 & 13444.1 $\pm$ 2613.2 & 1.9 $\pm$ 0.0 & 4.9 $\pm$ 0.0 & 0.18 $\pm$ 0.00 \\
fresh & 1522.4 $\pm$ 0.9 & 4370.6 $\pm$ 1.7 & 1619.1 $\pm$ 1.0 & 1643.1 $\pm$ 1.0 & 1412.80 $\pm$ 1.03 \\
switching & 1525.2 $\pm$ 20.3 & 12916.0 $\pm$ 1456.6 & 2291.7 $\pm$ 104.5 & 669.6 $\pm$ 13.1 & 1845.99 $\pm$ 84.56 \\
none & 76.1 $\pm$ 0.0 & 359.8 $\pm$ 0.3 & 1.5 $\pm$ 0.0 & 4.4 $\pm$ 0.0 & 0.18 $\pm$ 0.00 \\
\midrule
\multicolumn{6}{l}{Closed forms: $c_F T$ = 1551.7; residual (additive) $\sigma^2 T$ = 360.0; $\sigma^2 d$ = 0.18; $T$ = 4000, $M$ = 2000} \\
\bottomrule
\end{tabular}
}{\todo{run scripts/run\_theory.py}}
\caption{Cumulative regret $R(T)$ and $t\cdot\text{excess}_t$ late in the run (constant $\Rightarrow$ $1/t$ decay;
Theorem~\ref{thm:separation} predicts $\sigma^2d$). Mean $\pm$ s.d.\ over seeds.}
\label{tab:theory}
\end{table}

\subsection{Hierarchical planning under persistent surprises}
\label{sec:theory-hier}
A hierarchical planner chooses a route $r\in\{1,\dots,R\}$ (a subgoal sequence) and a low-level controller executes it.
Route $r$ succeeds iff the hidden disturbance lies in its feasible set, $\phi\in\mathcal F_r$; the bridge-like route
$r_0$ is feasible for every $\phi$ but costs more. Executing an infeasible route produces a \emph{surprise} (an observed
transition the nominal model does not predict) after $\tau_{\rm detect}$ steps, and an irreversible failure after
$\tau_{\rm irr}$ steps unless the planner aborts.

\begin{theorem}[Subgoal failure is permanent without hindsight]\label{thm:hier}
Let $\phi$ be persistent. A hierarchical planner whose subgoal choice is a fixed function of the prior (nominal model)
commits to $r^\star$ and re-selects $r^\star$ after every recoverable subgoal failure, so its main-goal failure
probability equals $\Pr(\phi\notin\mathcal F_{r^\star})$ \emph{for any number of retries}. A hierarchical planner whose
subgoal choice conditions on the observed surprise (posterior over $\phi$) and which aborts as soon as
$\Pr(\phi\notin\mathcal F_{r}\mid\text{history})$ exceeds a threshold fails only on the event
$\{\phi\notin\mathcal F_{r^\star}\}\cap\{\tau_{\rm detect}+\tau_{\rm abort}\ge\tau_{\rm irr}\}$.
Hence its advantage is $\Pr\big(\phi\notin\mathcal F_{r^\star},\ \tau_{\rm detect}+\tau_{\rm abort}<\tau_{\rm irr}\big)$,
which is positive iff the surprise is identifiable before the failure becomes irreversible.
If $\phi$ is redrawn independently at every attempt (fresh surprises), retries are independent and the
prior-based planner's failure probability after $n$ retries is $p^{\,n}$; the posterior carries no information about the
next draw, so hindsight gives no advantage (Theorem~\ref{thm:fresh}).
\end{theorem}
\begin{proof}[Proof sketch]
Under persistence the prior-based choice and the low-level outcome are deterministic functions of $\phi$ given the
start state, so every retry reproduces the same failure. The hindsight planner's abort rule triggers on the first
posterior update that makes $r^\star$ infeasible with high probability; it then re-plans over the remaining routes
and $r_0$ is always feasible. Under fresh draws the posterior predictive for the next attempt equals the prior.
\end{proof}
The theorem turns the qualitative claim ``hindsight makes hierarchical planning robust in deployment'' into two
testable conditions: \emph{persistence} of the surprise and \emph{detection before irreversibility}.
